% !TeX root = ../../Book1.tex
% 纲要标题：### 1.1 运算与结构
% 纲要位置：计划纲要.md，第 133 行。

\section{运算与结构}
\label{b1:sec:0101}

% 纲要内容（仅作写作计划，尚非教材正文）：
%
% * 集合上的一元、二元运算及其公理
% * 从整数加法、置换复合到矩阵乘法
% * 单位元、左右逆元的唯一性与消去
%

在整数中可以相加，在平面上可以先作一次旋转再作一次反射，在矩阵中可以相乘。这些计算所处理的对象不同，却都把给定的对象组合成同类对象。代数学把这种组合方式与对象本身一同研究：先确定运算在哪个集合中进行，再考察它满足哪些规律。只给出一个计算公式，还不一定定义了这个集合上的运算。

\ProseParagraphBreak
以下把\term[youxu-dui]{有序对}[ordered pair] \((a,b)\) 按两个位置分别比较，即 \((a,b)=(a',b')\) 当且仅当 \(a=a'\)、\(b=b'\)。给定集合 \(A,B\)，令
\[
A\times B=\bigl\{\,(a,b)\mid a\in A,\ b\in B\,\bigr\}
\]
称集合 \(A\times B\) 为 \(A,B\) 的\term[dikaer-ji]{笛卡尔积}[Cartesian product]。\footnote{笛卡尔（René Descartes，1596-1650），法国数学家与哲学家，以代数方法研究几何，开创解析几何。}
因此，对两个元素进行运算，可以表述为从有序对到运算结果的映射。

\subsection{一元、二元运算及其公理}
\label{b1:subsec:0101:01}

\begin{definition}\label{b1:def:operation}
给定集合 \(A\)，将映射 \(A\rightarrow A\) 称为 \(A\) 上的\term[yiyuan-yunsuan]{一元运算}[unary operation]，将映射 \(A\times A\rightarrow A\) 称为 \(A\) 上的\term[eryuan-yunsuan]{二元运算}[binary operation]。二元运算的值可写成 \(a*b\)。若对所有 \(a,b,c\in A\) 有
\[
(a*b)*c=a*(b*c),
\]
则称运算满足\term[jiehelv]{结合律}[associative law]；若对所有 \(a,b\in A\) 有 \(a*b=b*a\)，则称其满足\term[jiaohuanlv]{交换律}[commutative law]。
\end{definition}

“映射”要求每一对输入都有唯一确定、且属于 \(A\) 的输出。因此整数上的除法不是二元运算：除数可能为零，即使除数非零，结果也未必是整数。相反，整数减法确实是二元运算，只是不满足结合律。例如 \((3-2)-1=0\)，而 \(3-(2-1)=2\)。不满足某条公理与根本没有定义成运算，应当分开判断。

\ProseParagraphBreak
运算的公理是对任意元素都成立的等式或条件；计算几个例子只能发现反例，不能验证无限集合上的全称命题。对于有限集合，可以把二元运算列成表格；结合律要求相应等式对所有三元组成立，仅由运算表的对称性不能推出。表格关于主对角线对称表达的是交换律。

\ProseParagraphBreak
保持四个元素的次序，可以写出不同的括号表达式，例如
\[
((a*b)*c)*d,\qquad (a*b)*(c*d),\qquad a*(b*(c*d)).
\]
它们最后一次进行的运算，分别将前三个元素与第四个元素、前两个与后两个、第一个与后三个合在一起。一般地，一种完整的括号方式可以逐层构造：一个元素本身就是一个表达式；两个以上的元素则先分成两个非空连续块，分别给两块安排完整的括号，再把两块的结果相乘。因此，含 \(n\geq2\) 个元素的完整表达式，最外层总是把前 \(r\) 个元素的某种括号表达式与后 \(n-r\) 个元素的某种括号表达式相乘，其中 \(1\leq r<n\)。下面证明，结合律足以使这些不同安排给出相同的值。

\begin{proposition}[有限乘积的括号]\label{b1:prop:parentheses}
设 \(*\) 是集合 \(A\) 上满足结合律的二元运算，\(n\geq1\) 为整数，且 \(a_1,\ldots,a_n\in A\)。按原次序排列这 \(n\) 个元素时，不论如何加括号，运算结果都相同。
\end{proposition}

\begin{proof}
定义 \(P(a_1)=a_1\)，并递归定义
\[
P(a_1,\ldots,a_n)=P(a_1,\ldots,a_{n-1})*a_n\qquad(n\geq2).
\]
这个 \(P\) 表示从左端开始逐次相乘的值；要证的是每一种完整括号表达式的值都等于它。先说明两个各自从左端逐次相乘的连续块合在一起，仍得到整个序列的 \(P\)。具体地，对任意整数 \(r,s\geq1\) 及 \(a_1,\ldots,a_{r+s}\in A\)，有
\[
P(a_1,\ldots,a_r)*P(a_{r+1},\ldots,a_{r+s})
=P(a_1,\ldots,a_{r+s}).
\]
固定 \(r\)，对 \(s\) 归纳，并在每一步对任意相应长度的元素序列证明此式。当 \(s=1\) 时，这就是 \(P\) 的递归定义。若它对长度 \(s\) 成立，记 \(u=P(a_1,\ldots,a_r)\)、\(v=P(a_{r+1},\ldots,a_{r+s})\)，则结合律与归纳假设给出
\[
\begin{aligned}
u*(v*a_{r+s+1})&=(u*v)*a_{r+s+1}\\
&=P(a_1,\ldots,a_{r+s})*a_{r+s+1}\\
&=P(a_1,\ldots,a_{r+s+1}).
\end{aligned}
\]
这证明了合并两块的公式。再对总长度 \(n\) 作强归纳。\(n=1\) 时只有一个元素，结论显然。若 \(n\geq2\)，任意一种括号方式的最外层运算都把序列分成两个非空连续块，两块的长度均小于 \(n\)。由强归纳假设，两块的值分别等于相应的 \(P\)；合并公式再把结果化为 \(P(a_1,\ldots,a_n)\)。因此所有括号方式均给出同一个值。
\end{proof}

这里元素的次序始终没有改变。今后省略有限乘积中的括号，依据的是结合律。若两个相邻因子彼此交换，则可以把它们换序。例如取 \(a,b,c,d\in A\)，且 \(b*c=c*b\)，便有
\[
a*b*c*d=a*(b*c)*d=a*(c*b)*d=a*c*b*d.
\]
这一步用到的交换等式只有 \(b*c=c*b\)，对其余因子没有交换性要求；它保证这两个相邻因子的换序有效。若运算本身满足交换律，则可以任意重排有限乘积中的因子。

\subsection{整数、置换与矩阵上的运算}
\label{b1:subsec:0101:02}

整数集合 \(\mathbb Z\) 上，加法和乘法都满足结合律与交换律。取相反数 \(a\mapsto -a\) 是 \(\mathbb Z\) 上的一元运算；乘法的倒数则只有在扩大对象范围后才可能普遍存在，而且零仍然不能求倒数。这说明“逆运算”的可用范围也是结构的一部分。

\begin{definition}\label{b1:def:permutation}
给定集合 \(X\)，将 \(X\) 到自身的双射称为 \(X\) 的\term[zhihuan]{置换}[permutation]。记所有这样的双射构成的集合为 \(S_X\)。映射复合一律从右向左计算：\((fg)(x)=f\bigl(g(x)\bigr)\)。
\end{definition}

\symidx{\(S_X\)：集合 \(X\) 上的全部置换}

双射的复合仍是双射。对任意 \(f,g,h\) 和 \(x\in X\)，都有
\[
\bigl((fg)h\bigr)(x)=f\Bigl(g\bigl(h(x)\bigr)\Bigr)=\bigl(f(gh)\bigr)(x),
\]
故复合满足结合律。令 \(X=\{1,2,3\}\)，令 \(f\) 交换 \(1,2\) 而固定 \(3\)，令 \(g\) 交换 \(2,3\) 而固定 \(1\)。那么 \((fg)(1)=2\)，而 \((gf)(1)=3\)，所以复合一般不交换。这里两个映射在同一元素上的值不同，便足以说明它们不同。

\ProseParagraphBreak
固定大小的实方阵提供另一个例子。所有 \(n\) 阶实方阵对乘法封闭，并有
\[
\begin{aligned}
\bigl((AB)C\bigr)_{ij}
&=\sum_{k=1}^{n}\sum_{\ell=1}^{n}a_{i\ell}b_{\ell k}c_{kj}\\
&=\sum_{\ell=1}^{n}\sum_{k=1}^{n}a_{i\ell}b_{\ell k}c_{kj}
=\bigl(A(BC)\bigr)_{ij}.
\end{aligned}
\]
这里先按矩阵乘法的分量公式展开，再交换有限求和的次序，便得到同一个双重和。因此，矩阵乘法的结合律来自实数运算的结合律、分配律和有限求和的规则。取
\[
A=\begin{pmatrix}1&1\\0&1\end{pmatrix},\qquad
B=\begin{pmatrix}1&0\\1&1\end{pmatrix},
\]
则 \(AB=\begin{pmatrix}2&1\\1&1\end{pmatrix}\)，而
\(BA=\begin{pmatrix}1&1\\1&2\end{pmatrix}\)。即使两个矩阵都有逆，交换律也不会随之成立。

\subsection{单位元、逆元与消去}
\label{b1:subsec:0101:03}

\begin{definition}\label{b1:def:identity-inverse}
给定集合 \(A\) 上的二元运算 \(*\)。如果 \(e\in A\) 对每个 \(a\in A\) 都满足 \(e*a=a*e=a\)，则将 \(e\) 称为\term[danweiyuan]{单位元}[identity element]。假设单位元 \(e\) 存在，取 \(a,b\in A\)。若 \(b*a=e\)，则称 \(b\) 为 \(a\) 的左逆元；若 \(a*b=e\)，则称 \(b\) 为 \(a\) 的右逆元。若两式同时成立，则称 \(b\) 为 \(a\) 的双边\term[niyuan]{逆元}[inverse element]，简称逆元。左逆元和右逆元的定义都不排除另一侧等式成立。
\end{definition}

\begin{proposition}\label{b1:prop:identity-inverse-unique}
设 \(*\) 是集合 \(A\) 上的二元运算。这个运算的单位元至多有一个。如果 \(*\) 满足结合律并有单位元，那么任意 \(a\in A\) 的左逆元与右逆元只要同时存在就必相等；因而 \(a\) 的双边逆元若存在，则唯一。
\end{proposition}

\begin{proof}
若 \(e,f\) 都是单位元，则 \(e=e*f=f\)。若 \(b*a=e\) 且 \(a*c=e\)，则
\[
b=b*e=b*(a*c)=(b*a)*c=e*c=c.
\]
任意两个双边逆元分别充当左逆与右逆，即得唯一性。
\end{proof}

\begin{definition}\label{b1:def:cancellation}
给定集合 \(A\) 上的二元运算 \(*\)。若对任意 \(a,x,y\in A\)，\(a*x=a*y\) 都蕴含 \(x=y\)，则称该运算满足\term[zuo-xiaoqulv]{左消去律}[left cancellation law]；若对任意 \(a,x,y\in A\)，\(x*a=y*a\) 都蕴含 \(x=y\)，则称其满足\term[you-xiaoqulv]{右消去律}[right cancellation law]。同时满足左、右消去律时，称其满足\term[xiaoqulv]{消去律}[cancellation law]。
\end{definition}

单位元使空乘积有了自然的值，也使“撤销一次操作”可以写成等式。结合律则保证撤销操作能够穿过括号：当 \(a\) 有逆元时，\(a*x=a*y\) 两边左乘 \(a^{-1}\)，便得到 \(x=y\)。若只知道交换律，这个消去步骤没有根据；整数乘法中的 \(0\cdot 2=0\cdot 3\) 就不能消去零。

\ProseParagraphBreak
左右逆也确实可能不同步出现。令 \(\mathbb N=\{0,1,2,\ldots\}\)，取映射
\(s(n)=n+1\)、\(p(0)=0\)、\(p(n)=n-1\ (n\geq1)\)。对任意 \(n\in\mathbb N\)，有
\[
(p\circ s)(n)=p(n+1)=n,
\qquad (s\circ p)(0)=s(0)=1.
\]
所以 \(p\circ s=\operatorname{id}_{\mathbb N}\)，而 \(s\circ p\ne\operatorname{id}_{\mathbb N}\)。在 \(p\circ s\) 中，\(p\) 写在 \(s\) 的左边，因此称为 \(s\) 的左逆；按照从右向左的复合约定，实际先作用的却是 \(s\)，然后才由 \(p\) 将它撤销。右逆要求另一侧等式 \(s\circ p=\operatorname{id}_{\mathbb N}\) 成立，这里在 \(0\) 处已经失败。后面引入群时要求每个元素都有双边逆，正是为了排除这种不对称。
